% =========================
% Гэрийн даалгавар: Комплекс тоо (12 нүүрээс)
% =========================
\chapter*{Гэрийн даалгавар: Комплекс тоо}
\addcontentsline{toc}{chapter}{Гэрийн даалгавар: Комплекс тоо}
Дараах 20 бодлогыг бодож шалгуулна. Хүн бүрээс аль ч хамаагүй 5 бодлогыг бодуулж шалгана.
\section*{1.1 Суурь алгебр ба геометр чанарууд}
\begin{enumerate}[label=\textbf{1.\arabic*}, itemsep=8pt]

\item Дараах тэнцэтгэлүүдийг шалга.
\begin{enumerate}[label=\textbf{(\alph*)}]
\item $\big(\sqrt{2}-i\big)-i\big(1-\sqrt{2}i\big)=-2i$.
\item $(2-3i)(-2+i)=-1+8i$.
\end{enumerate}

\item Дараах илэрхийллийг хялбарчил:
\[
\frac{5i}{(1-i)(2-i)(3-i)}.
\]

\item Дараахыг батал:
\begin{enumerate}[label=\textbf{(\alph*)}]
\item $\Re(iz)=-\Im(z)$;
\item $\Im(iz)=\Re(z)$.
\end{enumerate}

\item Комплекс тоонуудын үржвэрийн \textbf{хаалт тавих} (associative) хууль хүчинтэйг батал:
\[
(z_1z_2)z_3=z_1(z_2z_3)\qquad(\forall z_1,z_2,z_3\in\mathbb C).
\]

\item Тооцоол:
\begin{enumerate}[label=\textbf{(\alph*)}]
\item $\displaystyle \frac{2+i}{2-i}$;
\item $(1-2i)^4$.
\end{enumerate}

\item Дараах дүрслэл өгөгдсөн байг:
\[
f:\mathbb{C}\to \mathbf{M}_{2,2} \,
, \; z=x+yi\ \mapsto\
\begin{pmatrix}
x & y\\
-y & x
\end{pmatrix}.
\]
Тэгвэл бүх $z_1,z_2\in\mathbb C$-ийн хувьд $f(z_1z_2)=f(z_1)f(z_2)$ болохыг батал.

\item Биномын томьёо ашиглан задла:
\begin{enumerate}[label=\textbf{(\alph*)}]
\item $(1+\sqrt{3}\,i)^{2011}$;
\item $(1+\sqrt{3}\,i)^{-2011}$.
\end{enumerate}

\item $z_1,z_2$ комплекс тоонууд ба $z_1\overline{z_2}$ нь тэг биш бодит бол $\exists r\in\mathbb{R},  z_1=r z_2$ болохыг батал.

\end{enumerate}

\section*{1.2 Модуль}
\begin{enumerate}[label=\textbf{2.\arabic*}, itemsep=8pt]

\item Бүх $z_1,z_2\in\mathbb C$-ийн хувьд
\[
|z_1-z_2|^2+|z_1+z_2|^2=2\big(|z_1|^2+|z_2|^2\big)
\]
болохыг батал.

\item Дараах тэнцэтгэл бишийг шалга:
\[
|\Re z|+|\Im z|\le \sqrt{2}\,|z|.
\]
\textit{Санамж:} Үүнийг $(|x|-|y|)^2\ge 0$.

\item Комплекс хавтгай дээрх муруйнуудыг \textbf{зурагла}:
\begin{enumerate}[label=\textbf{(\alph*)}]
\item $\Im(z)=-1$;
\item $|z-1|=|z+i|$;
\item $2|z|=|z-2|$.
\end{enumerate}

\item $|z|=R>1$ нөхцөлтэй бүх $z$-ийн хувьд дараах үнэлгээнүүдийг батал:
\[
\frac{R^4-R}{R^2+R+1}\le
\left|\frac{z^4+iz}{z^2+z+1}\right|
\le
\frac{R^4+R}{(R-1)^2}.
\]

\item $z\ne 0$ үед
\[
|\log(z)|\le |\ln|z||+\pi
\]
болохыг батал. (Энд $\log$ нь үндсэн логарифм.)

\end{enumerate}

\section*{1.3 Экспоненциал ба туйлын хэлбэр, комплекс язгуурууд}
\begin{enumerate}[label=\textbf{3.\arabic*}, itemsep=8pt]

\item Дараахыг $x+iy$ хэлбэрт (энд $x,y\in\mathbb R$) хувирга:
\begin{enumerate}[label=\textbf{(\alph*)}]
\item $\displaystyle \frac{i}{1-i}+\frac{1-i}{i}$;
\item $-8i$ тооны бүх \textbf{гуравдугаар зэрийн} язгуурууд;
\item $\left(\frac{i+1}{\sqrt2}\right)^{2026}$.
\end{enumerate}

\item Дараах тоонуудын \textbf{ аргумент} (principal argument) болон \textbf{илтгэгч хэлбэр}-ийг ол:
\begin{enumerate}[label=\textbf{(\alph*)}]
\item $\displaystyle z=\frac{i}{1+i}$;
\item $z=\sqrt3+i$;
\item $z=2-i$.
\end{enumerate}

\item Дараах тэгшитгэлүүдийн бүх комплекс шийдийг ол:
\begin{enumerate}[label=\textbf{(\alph*)}]
\item $z^6=-9$;
\item $z^2+2z+(1-i)=0$.
\end{enumerate}

\item $z^4+16=0$ тэгшитгэлийн 4 язгуурыг ол.
Дараа нь эдгээрийг ашиглан $z^4+16$-г \textbf{бодит коэффициенттэй} хоёр квадрат үржвэрт задла:
\[
z^4+16=(z^2+az+b)(z^2+cz+d)\quad (a,b,c,d\in\mathbb R).
\]

\item Дараахыг хий:
\begin{enumerate}[label=\textbf{(\alph*)}]
\item Илтгэгч хэлбэр ашиглан
\[
(1+\sqrt3\,i)^{2011},\qquad (1+\sqrt3\,i)^{-2011}
\]
-ийг ол.
\item Дараах хоёр тэнцэтгэлүүдийг батал:
\[
\sum_{m=0}^{1005}\binom{2011}{2m}(-3)^m=2^{2010},
\qquad
\sum_{m=0}^{1005}\binom{2011}{2m+1}(-3)^m=2^{2010}.
\]
\end{enumerate}

\item Эхлээд дараах нийлбэрийн томьёог батал:
\[
1+z+z^2+\cdots+z^n=\frac{1-z^{n+1}}{1-z}\qquad (z\ne 1),
\]
дараа нь үүнийг ашиглан \textbf{Лагранжийн тригонометрийн адилтгал}-ыг гарга:
\[
1+\cos q+\cos 2q+\cdots+\cos nq
=
\frac12+\frac{\sin\left(\frac{(2n+1)q}{2}\right)}{2\sin\left(\frac{q}{2}\right)}
\qquad (0<q<2\pi).
\]

\item Комплекс тоо ашиглан \textbf{косинусын теорем}-ыг батал:
Гурвалжин $\triangle ABC$-д $|BC|=a$, $|CA|=b$, $|AB|=c$, $\angle BCA=q$ бол
\[
a^2+b^2-2ab\cos q=c^2.
\]
\textit{Санамж:} $C$-г тэг дээр байрлуулаад $B=z_1$, $A=z_2$ гэж авч,
$z_1\overline{z_2}+z_2\overline{z_1}=2|z_1||z_2|\cos q$ гэдгийг харуул.

\end{enumerate}